%% ============================================================
%%  REMERCIEMENTS
%% ============================================================
\chapter*{Remerciements}
\addcontentsline{toc}{chapter}{Remerciements}
\thispagestyle{fancy}

Nous tenons à exprimer notre profonde gratitude à toutes les personnes qui ont contribué,
de près ou de loin, à la réalisation de ce projet.

\begin{tcolorbox}[infobox, title={\faHeart\quad Encadrement}]
Nos remerciements les plus sincères vont à \textbf{Monsieur Mohamed Lachgar}, notre encadrant
de projet, pour ses précieux conseils, sa disponibilité constante et son suivi rigoureux
tout au long de ce travail. Sa vision claire des enjeux de la sécurité mobile et son
expertise technique ont été une source d'inspiration et de motivation irremplaçable.
\end{tcolorbox}

\vspace{0.5cm}

\begin{tcolorbox}[infobox, title={\faUniversity\quad Corps enseignant}]
Nous adressons également nos remerciements à l'ensemble du corps enseignant de notre
établissement, pour la qualité de la formation dispensée et les connaissances transmises
qui ont rendu ce projet possible.
\end{tcolorbox}

\vspace{0.5cm}

\begin{tcolorbox}[successbox, title={\faUsers\quad Communautés Open Source}]
Ce projet repose sur des outils et standards libres de grande qualité. Nous remercions
chaleureusement :
\begin{itemize}[leftmargin=1.5em]
  \item La \textbf{fondation OWASP} pour le MASVS v2 et le MASTG,
  \item L'équipe \textbf{MobSF} pour leur remarquable outil d'analyse statique,
  \item Les contributeurs de \textbf{FastAPI}, \textbf{Next.js}, \textbf{Celery} et \textbf{Docker},
  \item L'équipe \textbf{Ollama} pour la démocratisation des LLM locaux.
\end{itemize}
\end{tcolorbox}

\vspace{0.5cm}

Enfin, nous remercions nos familles pour leur patience infinie, et nos amis pour leur
soutien moral pendant les moments les plus intenses de ce projet.

\cleardoublepage
